Как только базовая модель платежа выстроена, возникает следующий
вопрос: почему этой системе вообще стоит доверять.
Покупатель прикладывает карту к~терминалу, на экране загорается сумма,
авторизационный запрос уходит к~эмитенту, и~вся цепочка должна за
доли секунды решить, что перед ней не клон, не подменённые данные и~не
имитация согласия держателя. Знание того, кто несёт
денежное обязательство, кто посылает сообщение и кто меняет запись
в~своём учёте, ещё не делает операцию возможной.

Часть~I описала внешний поток между банками, сетями и продавцом.
Теперь в~фокусе внутренние проверки системы: как
терминал и~карта договариваются о~правилах
операции, как эмитент получает основания сказать \enquote{да} или
\enquote{нет}, как токен заменяет PAN, не разрушая жизненный цикл
транзакции, как криптография заменяет статический идентификатор
динамическими доказательствами подлинности и как защита данных в~платеже
распределяется между несколькими устройствами и стандартами.

Эта часть собрана вокруг инвариантов, которые повторяются в~разных
формах независимо от конкретной сети. EMV отвечает на вопрос, как
карта доказывает подлинность в~точке приёма. Бесконтактные платежи
показывают, как та же логика переносится в~сценарий, где скорость
и~удобство нельзя купить ценой доверия к~операции. Токенизация
разбирает другой тип проблемы: как дать участникам возможность
обрабатывать платёж, не превращая PAN в~универсальный ключ доступа к~счёту.
Криптография даёт этим сюжетам общий язык секретов, ключей,
криптограмм и проверок. 3-D Secure переносит проверку в~электронную
коммерцию, где карта и~терминал не встречаются физически.
PCI DSS отвечает на вопрос, какие данные допустимо удерживать
внутри системы и~какую цену приходится платить за~их присутствие.

Часть~II отвечает на вопрос, как платёж совмещает удобство
и~проверяемость~--- два требования, которые тянут в~разные стороны.
Чип не должен делать оплату медленной, токенизация~--- ломать
регулярные списания, 3-D Secure~--- превращать страницу оплаты
в~испытание, PCI DSS~--- оставаться внешним аудитом, отделённым
от архитектуры продукта.

Часть остаётся на уровне общих механизмов: офлайн- и~онлайн-проверка,
криптографическое доказательство, замена чувствительного
идентификатора, удалённое подтверждение держателя, сокращение зоны
хранения данных. Конкретные правила \textit{fallback} и пороги
сертификации каждой сети~--- в~следующей части.

В~части~\ref{part:payment-systems} те же механизмы начнут жить
в~конкретных системах: Visa, Mastercard, СБП, открытом банкинге.
Один принцип безопасности влечёт там разные операционные и коммерческие
последствия в~зависимости от свода правил и расчётной модели.
